% GENERATED FILE. Edit canonical lessons, then run npm run generate:books.
% canonical-source-hash: fnv1a64:eedcdf54eca7687f
% canonical-lessons: RU-C169-prosba, RU-C169-pomoshch, RU-C169-mneniye, RU-C169-pravda, RU-C169-lozh

\chapter{Request, Help, Opinion, Truth, Lie}
\label{ch:ru-request-help-opinion-truth-lie}
\hlchaptermodality{169}

\begin{chapteropening}
\textbf{By the end of this chapter:} \emph{I can say a request, help, an opinion, the truth and a lie.}

Complete the last lesson of chapter 169: I can say a request, help, an opinion, the truth and a lie.
\end{chapteropening}

\section[\texorpdfstring{\ru{просьба} (prósba)}{просьба (prósba)}]{\ru{просьба} (prósba) — a request}

\label{lesson:RU-C169-prosba}



\begin{quote}
Before the new one: say the Russian for a relative, then the Russian for a pensioner.
\end{quote}

\subsection*{You'll want to know: \ru{просьба}}
\textbf{\ru{просьба}} (\emph{prósba}) — \textquotedblleft{}a request\textquotedblright{}.

\begin{grammarlens}[title={What you've built}]
One more word for this chapter.
\end{grammarlens}

\subsection*{Your turn}
\begin{itemize}
\raggedright
  \item \emph{Say it:} \emph{prósba}
  \item \emph{Say it:} \emph{prósba}, once more
  \item \emph{Say it:} say \emph{prósba}
  \item \emph{Recall:} say the Russian for a relative, then the Russian for a pensioner, then say \emph{prósba} again
\end{itemize}

\begin{tcolorbox}[breakable,colback=teal!4,colframe=teal!35!black,title={Before you move on}]
What does \ru{просьба} mean? (\textquotedblleft{}A request\textquotedblright{}.) Say it once more.
\end{tcolorbox}
\section[\texorpdfstring{\ru{помощь} (pómoshch)}{помощь (pómoshch)}]{\ru{помощь} (pómoshch) — help}

\label{lesson:RU-C169-pomoshch}



\begin{quote}
Before the new one: say the Russian for a pensioner, then the Russian for a request.
\end{quote}

\subsection*{You'll want to know: \ru{помощь}}
\textbf{\ru{помощь}} (\emph{pómoshch}) — \textquotedblleft{}help\textquotedblright{}.

\begin{grammarlens}[title={What you've built}]
One more word for this chapter.
\end{grammarlens}

\subsection*{Your turn}
\begin{itemize}
\raggedright
  \item \emph{Say it:} \emph{pómoshch}
  \item \emph{Say it:} \emph{pómoshch}, once more
  \item \emph{Say it:} say \emph{pómoshch}
  \item \emph{Recall:} say the Russian for a pensioner, then the Russian for a request, then say \emph{pómoshch} again
\end{itemize}

\begin{tcolorbox}[breakable,colback=teal!4,colframe=teal!35!black,title={Before you move on}]
What does \ru{помощь} mean? (\textquotedblleft{}Help\textquotedblright{}.) Say it once more.
\end{tcolorbox}
\section[\texorpdfstring{\ru{мнение} (mnéniye)}{мнение (mnéniye)}]{\ru{мнение} (mnéniye) — an opinion}

\label{lesson:RU-C169-mneniye}



\begin{quote}
Before the new one: say the Russian for a request, then the Russian for help.
\end{quote}

\subsection*{You'll want to know: \ru{мнение}}
\textbf{\ru{мнение}} (\emph{mnéniye}) — \textquotedblleft{}an opinion\textquotedblright{}.

\begin{grammarlens}[title={What you've built}]
One more word for this chapter.
\end{grammarlens}

\subsection*{Your turn}
\begin{itemize}
\raggedright
  \item \emph{Say it:} \emph{mnéniye}
  \item \emph{Say it:} \emph{mnéniye}, once more
  \item \emph{Say it:} say \emph{mnéniye}
  \item \emph{Recall:} say the Russian for a request, then the Russian for help, then say \emph{mnéniye} again
\end{itemize}

\begin{tcolorbox}[breakable,colback=teal!4,colframe=teal!35!black,title={Before you move on}]
What does \ru{мнение} mean? (\textquotedblleft{}An opinion\textquotedblright{}.) Say it once more.
\end{tcolorbox}
\section[\texorpdfstring{\ru{правда} (právda)}{правда (právda)}]{\ru{правда} (právda) — the truth}

\label{lesson:RU-C169-pravda}



\begin{quote}
Before the new one: say the Russian for help, then the Russian for an opinion.
\end{quote}

\subsection*{You'll want to know: \ru{правда}}
\textbf{\ru{правда}} (\emph{právda}) — \textquotedblleft{}the truth\textquotedblright{}.

\begin{grammarlens}[title={What you've built}]
One more word for this chapter.
\end{grammarlens}

\subsection*{Your turn}
\begin{itemize}
\raggedright
  \item \emph{Say it:} \emph{právda}
  \item \emph{Say it:} \emph{právda}, once more
  \item \emph{Say it:} say \emph{právda}
  \item \emph{Recall:} say the Russian for help, then the Russian for an opinion, then say \emph{právda} again
\end{itemize}

\begin{tcolorbox}[breakable,colback=teal!4,colframe=teal!35!black,title={Before you move on}]
What does \ru{правда} mean? (\textquotedblleft{}The truth\textquotedblright{}.) Say it once more.
\end{tcolorbox}
\section[\texorpdfstring{\ru{ложь} (lozh)}{ложь (lozh)}]{\ru{ложь} (lozh) — a lie}

\label{lesson:RU-C169-lozh}



\begin{quote}
Before the new one: say the Russian for an opinion, then the Russian for the truth.
\end{quote}

\subsection*{You'll want to know: \ru{ложь}}
\textbf{\ru{ложь}} (\emph{lozh}) — \textquotedblleft{}a lie\textquotedblright{}.

\begin{grammarlens}[title={What you've built}]
One more word for this chapter.
\end{grammarlens}

\subsection*{Your turn}
\begin{itemize}
\raggedright
  \item \emph{Say it:} \emph{lozh}
  \item \emph{Say it:} \emph{lozh}, once more
  \item \emph{Say it:} say \emph{lozh}
  \item \emph{Recall:} say the Russian for an opinion, then the Russian for the truth, then say \emph{lozh} again
\end{itemize}

\begin{tcolorbox}[breakable,colback=teal!4,colframe=teal!35!black,title={Before you move on}]
What does \ru{ложь} mean? (\textquotedblleft{}A lie\textquotedblright{}.) Say it once more.
\end{tcolorbox}
